% Corrected French for printed p.132.
% Corrections: avec seule flèche -> avec pour seule flèche; formé -> formée.
\SourcePage{132}

\LayerTarget{tohoku:example:1.7:b:p2}
\noindent relations de commutation, on obtient un produit de catégories du
type envisagé dans 1.

\LayerTarget{tohoku:example:1.7:c}
c) Prenons $I$ réduit à deux éléments, $a$ et $b$, avec une seule
flèche d'origine $a$ et extrémité $b$ : on trouve la catégorie des morphismes
$u:A\to B$ entre objets de $\mathbf{C}$. L'introduction de relations de
commutation reviendrait à se borner aux $A,B,u$ annulés par certains entiers
précisés.

\LayerTarget{tohoku:example:1.7:d}
d) \textit{Catégories de foncteurs.} Soit $\mathbf{C}'$ une autre catégorie,
supposons qu'elle soit un ensemble. Alors les foncteurs covariants de $\mathbf{C}'$ dans
$\mathbf{C}$ forment une catégorie, en prenant pour morphismes les morphismes
fonctoriels (cf. 1, 1). Cette catégorie peut s'interpréter comme une catégorie
$\mathbf{C}^\Sigma$, où l'on prend $I=\mathbf{C}'$, les flèches d'origine $A'$ et d'extrémité
$B'$ étant par définition les éléments de $\mathrm{Hom}(A',B')$, et les
relations de commutation étant celles qui expriment les deux axiomes d'un
foncteur. Si $\mathbf{C}'$ est de plus une catégorie additive, les foncteurs
\textit{additifs} de $\mathbf{C}'$ dans $\mathbf{C}$ peuvent aussi s'interpréter comme une
catégorie $\mathbf{C}^\Sigma$ (on rajoute les relations de commutation qu'il faut).

\LayerTarget{tohoku:example:1.7:e}
e) \textit{Complexes à valeurs dans $\mathbf{C}$.} $I=Z$ (ensemble
des entiers), l'ensemble des flèches étant $(d_n)_{n\in Z}$ où $d_n$
est d'origine $n$ et d'extrémité $n+1$, les relations de commutation étant
$d_{n+1}d_n=0$. On peut encore rajouter des relations de la forme $e_n=0$ si
on veut se borner aux complexes à degrés positifs, ou à degrés négatifs. On
obtient de façon analogue les bicomplexes etc.

\LayerTarget{tohoku:example:1.7:f}
f) \textit{La catégorie $\mathbf{C}^G$, ($G$ un groupe).} Soit $G$ un groupe,
$\mathbf{C}$ une catégorie (non nécessairement additive). On appelle \textit{objet avec
groupe d'opérateurs $G$} dans $\mathbf{C}$, un couple $(A,r)$ formé d'un objet $A\in \mathbf{C}$
et d'une représentation $r$ de $G$ dans le groupe des automorphismes de $A$.
Si $(A',r')$ est un deuxième tel couple, on appelle morphisme du premier dans
le second un morphisme de $A$ dans $A'$ qui permute aux opérations de $G$. La
classe $\mathbf{C}^G$ des objets dans $\mathbf{C}$ avec groupe d'opérateurs $G$ devient ainsi une
catégorie. On peut l'interpréter comme une classe $\mathbf{C}^\Sigma$, où on prend pour
$\Sigma=\Sigma(G)$ le schéma avec relations suivant : l'ensemble des sommets
est réduit à un élément $i_0$, l'ensemble des flèches est $G$, les relations de
commutation sont $(s)(t)=(st)$ (où le premier membre désigne une flèche
composée) et $(e)=e_{i_0}$ (où $e$ désigne l'élément unité de $G$). En
particulier, si $\mathbf{C}$ est une catégorie additive, il en est de même de $\mathbf{C}^G$ ;
dans ce cas, notre construction est contenue dans celle qui suit, (grâce à la
considération de l'algèbre du groupe $G$).

\LayerTarget{tohoku:example:1.7:g}
g) \textit{La catégorie $\mathbf{C}^U$ ($U$ un anneau avec unité).} On
considère la catégorie additive formée des couples $(A,r)$ d'un objet $A$ de
$\mathbf{C}$ et d'une représentation unitaire de $U$ dans l'anneau
$\mathrm{Hom}(A,A)$, les morphismes dans cette catégorie étant définis de
façon évidente. Elle s'interprète encore comme ci-dessus comme une catégorie
$\mathbf{C}^{\Sigma(U)}$, où $\Sigma(U)$ est le schéma avec relations ayant un seul
sommet, $U$ comme ensemble de flèches, et des relations de commutation que nous
n'écrirons pas.

\LayerTarget{tohoku:example:1.7:h:p1}
h) \textit{Systèmes inductifs et systèmes projectifs.} On prend
comme ensemble de sommets un ensemble \textit{préordonné} $I$, comme flèches
les couples $(i,j)$ de sommets avec $i\leqq j$, l'origine et l'extrémité de
$(i,j)$ étant respectivement $i$ et $j$. Les relations de commutation sont
$(i,j)(j,k)=(i,k)$ et $(i,i)=e_i$. Les

\TohokuContinuation{133}
